%%%%%%%%%%%%%%%%%%%%%%%%%%%%%%%%%%%%%%%%%%%%%%%%%%%%%%%%%%%%%%%%%%%%%%%%%%%
%%% Selçuk Üniversitesi Fen Bilimleri Enstitüsü LaTeX Tez Yazım Şablonu											  
%%% Oluşturan: Arş. Gör. Metin Turgay 		                                   
%%% Son düzenleme ve yükleme:  Aralık 2023			       
%%% En son FBE tez yazım dökümanına göre düzenlenmiştir.            	  

%%% Bu döküman Selçuk Eylül Üniversitesi Fen Bilimleri Enstitüsünün Aralık 
%%% 2023 kurallarına göre düzenlenniştir. Enstütü yazım kuralları sürekli 
%%% güncellendiği için, değişiklileri takip etmek yazarın 
%%% sorumluluğundadır.

%%% Bu şablonu derleyebilmek için 
%%% "Options -> Build -> Default Complier"
%%% seçeneğini "xelatex" olarak güncelleyiniz.

%% Aşağıdaki komutu değiştirmeyiniz!
\documentclass[12pt,oneside]{report}    

%% Selçuk Üniversitesi FBE Tez yazım şablonu paketini yüklüyoruz.
%% Aşağıdaki iki komutu silmeyiniz!
\usepackage{selcuk}             
\emergencystretch=1em


%--------------------------------------------------------------------------
% TEZ BİLGİLERİ
%--------------------------------------------------------------------------
%
% AŞAĞIDAKİ BİLGİLERİ DİKKATLİ ŞEKİLDE DOLDURMANIZ TEZ İÇERİSİNDE OTOMATİK 
% OLARAK DOLDURULAN KISIMLAR İÇİN ÇOK ÖNEMLİDİR!
%
%--------------------------------------------------------------------------

\renewcommand{\thesisauthor}{Öğrencinin Adı SOYADI}    %% Yazar adı soyadı.
\renewcommand{\thesisdepartment}{Matematik}   %% Bölüm.
\renewcommand{\thesismonth}{Mezuniyet Ayı}    %% Mezuniyet ayı.
\renewcommand{\thesisyear}{Mezuniyet Yılı}    %% Mezuniyet yılı.
\renewcommand{\thesisdefense}{Savunma Tarihi} %% Savunma tarihi. ../../.... formatında girilmelidir.
\renewcommand{\thesisdegree}{YÜKSEK LİSANS/DOKTORA} %% Birini siliniz.
\renewcommand{\thesistitle}{TEZ BAŞLIĞININ İNGİLİZCE’SİNİ BURAYA YAZINIZ}  
\renewcommand{\thesistitleturkish}{TEZ BAŞLIĞINI BURAYA BÜYÜK HARFLER İLE YAZINIZ}
\renewcommand{\thesissupervisor}{Danışman Adı SOYADI}  %% Danışman Adı
\renewcommand{\thesisdean}{Prof. Dr. Ömer Faruk YÜKSEL}
\renewcommand{\thesisauthoraddress}{KONYA}
\renewcommand{\bibname}{KAYNAKLAR}

%--------------------------------------------------------------------------
%	TEZ icerisindeki Bilgiler
%--------------------------------------------------------------------------

\begin{document}
        %%% Başlık sayfasını oluştur.
	\thesistitlepage     
	
        %%% Tez Kabul sayfasını oluştur.
	\begin{tezkabul}
	
        \thispagestyle{empty}
		
		\textbf{\thesisauthor} \, taraf\i ndan haz\i rlanan ``\textbf{\thesistitleturkish}" adlı tez çal\i şması  {\thesisdefense} \, tarihinde aşa\u{g}\i daki jüri taraf\i ndan oy	birli\u{g}i ile Selçuk Üniversitesi Fen Bilimleri Enstitüsü {\thesisdepartment} \, Anabilim Dal\i 'nda {\thesisdegree} TEZ\.I olarak kabul edilmiştir.
		
		\vspace*{3.5cm} 
		
		\noindent\textbf{Jüri Üyeleri} \hfill  \textbf{\.Imza \hspace{1.2cm}} \bigskip
		\\
		\noindent\textbf{Başkan}\\
		\textbf{Unvanı Adı SOYADI} \hfill  \textbf{\rule{0.2\textwidth}{.4pt}} \medskip \\
		
		\noindent\textbf{Dan\i şman} \\
		\noindent\textbf{\thesissupervisor} \hfill \hfill
		\textbf{\rule{0.2\textwidth}{.4pt}} \medskip \\
		
		\noindent\textbf{Üye} \\
		\noindent\textbf{Unvanı Adı SOYADI} \hfill \hfill \textbf{\rule{0.2\textwidth}{.4pt}}
		\medskip \\
		
		%% Yüksek Lisans tezi hazırlıyorsanız buradan itibaren iki juri üyesi satırlarını siliniz. 
		\noindent\textbf{Üye} \\
		\noindent\textbf{Unvanı Adı SOYADI} \hfill \hfill \textbf{\rule{0.2\textwidth}{.4pt}}
		\medskip \\
		
		\noindent\textbf{Üye} \\
		\noindent\textbf{Unvanı Adı SOYADI} \hfill \hfill \textbf{\rule{0.2\textwidth}{.4pt}}
		\medskip \\
		
		
		Yukar\i daki sonucu onaylar\i m.
		
		\begin{flushright} 
			\thesisdean \\
			FBE Müdürü
			
		\end{flushright} 	
		\vspace{3cm}
		Bu tez çal\i şması  ... taraf\i ndan ... nolu proje ile desteklenmiştir.  %% Proje yok ise yukarıdaki satırı siliniz.
	\end{tezkabul}
	
	%%% Bildirim sayfasını oluştur.
	\begin{bildirim}
 
        \thispagestyle{empty}
		Bu tezdeki bütün bilgilerin etik davranış ve akademik kurallar çerçevesinde elde edildiğini ve tez yazım kurallarına uygun olarak hazırlanan bu çalışmada bana ait olmayan her türlü ifade ve bilginin kaynağına eksiksiz atıf yapıldığını bildiririm.
		
		\bigskip
		\bigskip
		
		\begin{center}
			\bfseries\MakeUppercase DECLARATION PAGE
		\end{center}
		
		I hereby declare that all information in this document has been obtained and presented in accordance with academic rules and ethical conduct. I also declare that, as required by these rules and conduct, I have fully cited and referenced all material and results that are not original to this work.
		\bigskip
		\vspace{0.5cm}
		
		\hfill
		\begin{tabular}{@{}l@{}}
			İmza % Bu yazıyı siliniz!
			\\
			\thesisauthor\\ % Tez bilgileri kısmını doldurduysanız burayı otomatik olarak güncelleyecektir!
			Tarih: ../../20..
		\end{tabular}
		
	\end{bildirim}
	
	%%% Özet sayfasını oluştur.
	
	\begin{thesisozet} 
		
		\begin{center}
			\textbf{\MakeUppercase \thesisdegree \,TEZ\.I} \bigskip\\
			\textbf{\MakeUppercase \thesistitleturkish} \bigskip \bigskip\\
			\textbf{\thesisauthor} \bigskip \\
			\textbf{Sel\c{c}uk \"{U}niversitesi Fen Bilimleri Enstit\"{u}s\"{u}}\\
			\textbf{\thesisdepartment \,Anabilim Dal\i } \bigskip \\
			\textbf{Dan\i\c{s}man:} \textbf{\thesissupervisor} \bigskip\\
			\textbf{\thesisyear, ... Sayfa } \bigskip\\ 
			\bigskip
			\textbf{Jüri} \\
			\textbf{Danışmanın Unvanı Adı SOYADI} \\%Danışman
			\textbf{Diğer Üyenin Unvanı Adı SOYADI} \\
			\textbf{Diğer Üyenin Unvanı Adı SOYADI} \\
			\textbf{Diğer Üyenin Unvanı Adı SOYADI} \\
			\textbf{Diğer Üyenin Unvanı Adı SOYADI} \\
									
		\end{center} \par
		\bigskip \bigskip
		\footnotesize Özet metnini yazmaya buradan başlayınız Özet metnini yazmaya buradan başlayınız Özet metnini yazmaya buradan başlayınız Özet metnini yazmaya buradan başlayınız Özet metnini yazmaya buradan başlayınız Özet metnini yazmaya buradan başlayınızÖzet metnini yazmaya buradan başlayınız Özet metnini yazmaya buradan başlayınız Özet metnini yazmaya buradan başlayınız 
		
		\footnotesize \textbf{Anahtar Kelimeler:}  4-8 adet anahtar kelime yazınız. Alfabetik sırada olmalı
	\end{thesisozet} 
	
	%%% Abstract sayfasını oluştur.
	
	\begin{thesisabstract}
		\begin{center}	
			\textbf{\MakeUppercase MS/Ph.D THESIS} \bigskip \\
			\textbf{\MakeUppercase \thesistitle}\bigskip \bigskip\\
			\textbf{\thesisauthor} \bigskip\\
			\textbf{THE GRADUATE SCHOOL OF  NATURAL AND APPLIED SCIENCE OF SELÇUK UNIVERSITY}\\
			\textbf{THE DEGREE OF  MASTER OF SCIENCE / DOCTOR OF PHILOSOPHY IN MATHEMATICS} % Burayı bölümünüz ile güncelleyiniz!
                \bigskip\\
			\textbf{Advisor: Danışman Adı SOYADI}  % Danışmanın ingilizce ünvanı ile beraber doldurunuz!
   \bigskip\\
			\textbf{Year, ... Pages}\bigskip\\
			\bigskip
			\textbf{Jury} \\
			\textbf{Danışman Title Adı SOYADI} \\ %Danışman
			\textbf{Diğer Üyenin Title Adı SOYADI} \\
			\textbf{Diğer Üyenin Title Adı SOYADI} \\
			\textbf{Diğer Üyenin Title Adı SOYADI} \\
			\textbf{Diğer Üyenin Title Adı SOYADI} \\
		\end{center} \par
		\bigskip \bigskip
		\footnotesize Türkçe özet metninin İngilizce’sini yazmaya buradan başlayınız. Türkçe özet metninin İngilizce’sini yazmaya buradan başlayınız. Türkçe özet metninin İngilizce’sini yazmaya buradan başlayınız. Türkçe özet metninin İngilizce’sini yazmaya buradan başlayınız. Türkçe özet metninin İngilizce’sini yazmaya buradan başlayınız. Türkçe özet metninin İngilizce’sini yazmaya buradan başlayınız. Türkçe özet metninin İngilizce’sini yazmaya buradan başlayınız.  
		
		\footnotesize \textbf{Keywords:}  Türkçe özetteki anahtar kelimelerin İngilizce’sini yazınız. Alfabetik sırada olmalı
	\end{thesisabstract} 
	
	%%% Önsöz sayfasını oluştur.
	
	\begin{onsoz} 
		Önsöz metnini yazım kılavuzuna uygun olarak yazmaya buradan başlayınız.
		\begin{flushright} 
			\thesisauthor \\
			\thesisauthoraddress-\thesisyear
		\end{flushright}
	\end{onsoz}
	

	\begin{spacing}{1}
	\tableofcontents                   %% İçindekiler kısmını oluşturur.
	\end{spacing}
	
	% ŞEKİL VE TABLO LİSTESİ GEREKLİYSE AŞAĞIDAKİ KOMUTLARDAN % İFADESİNİ SİLİNİZ!
	%
	%\renewcommand{\listfigurename}{ŞEKİL LİSTESİ}
	%\listoffigures   
	%
	%\renewcommand{\listtablename}{TABLO LİSTESİ}	
	%\listoftables   
	%
	%%%
	%%% Generate Simgeler ve K\i saltmalar page.
	%%%
	\begin{simge} 
		\textbf{Simgeler}
		
		Simgeleri yazmaya buradan başlayınız ve yazım kılavuzunda belirtildiği şekilde düzenleyiniz. 
					
		\vspace{0.5cm}
		
		\textbf{K\i saltmalar}
		
		Kısaltmaları yazmaya buradan başlayınız ve yazım kılavuzunda belirtildiği şekilde düzenleyiniz.
		
	\end{simge}

        %%% Tez sayfasını oluştur.
	
	\rhead{\thepage}
	\fancyfoot{}
	\renewcommand{\baselinestretch}{1.5}
	
	%%%
	%%% BURADAN TEZİNİZİN İÇERİĞİNİ YAZMAYA BAŞLAYINIZ
	%%%
	\chapter{GİRİŞ}
	
	Giriş bölümünü yazmaya buradan başlayınız. 
	
	\section{Birinci Bölüm İkinci Derece Başlık}
	\begin{theorem}\label{teorem1}
		Örnek teorem ifadesi
	\end{theorem}
	
	\begin{proof}
		Örnek ispat
	\end{proof}

        Denklem yazımlarına örnek olması için aşağıdaki teorem ispat edilmiştir.
        \begin{theorem}
            ${\left( {f\left( x \right) \pm g\left( x \right)} \right)^\prime } = f'\left( x \right) \pm g'\left( x \right)$ ifadesinin doğruluğunu gösteriniz.
        \end{theorem}
        \begin{proof}
            Bu ifade, türevin tanımını kullanarak kolayca kanıtlanabilir. İki fonksiyonun toplamından başlayacağız. İlk olarak, türevin tanımına bu toplamı yerine koyup paydada biraz değişiklik yapacağız.
            \begin{align*} % denklemlere numara verilmeyecekse align* kullanılmalıdır.
                (f(x)+g(x))^{\prime} & =\lim _{h \rightarrow 0} \frac{f(x+h)+g(x+h)-(f(x)+g(x))}{h} \\
                & =\lim _{h \rightarrow 0} \frac{f(x+h)-f(x)+g(x+h)-g(x)}{h}
            \end{align*}
            Şimdi, kesiri iki parçaya bölelim ve bir toplamın limitinin toplamının limiti olduğunu hatırlayalım. Bu gerçeği kullanarak, her iki fonksiyon için türev tanımına ulaştığımızı görebiliriz.
            \begin{align}
                (f(x)+g(x))^{\prime} & =\lim _{h \rightarrow 0} \frac{f(x+h)-f(x)}{h}+\lim _{h \rightarrow 0} \frac{g(x+h)-g(x)}{h} \\
                & =f^{\prime}(x)+g^{\prime}(x)
            \end{align}
            İki fonksiyonun farkının kanıtı neredeyse aynı olduğu için, bunu burada herhangi bir açıklama yapmadan vereceğiz.
            \begin{align}
                (f(x)-g(x))^{\prime} & =\lim _{h \rightarrow 0} \frac{f(x+h)-g(x+h)-(f(x)-g(x))}{h} \\
                & =\lim _{h \rightarrow 0} \frac{f(x+h)-f(x)-(g(x+h)-g(x))}{h} \\
                & =\lim _{h \rightarrow 0} \frac{f(x+h)-f(x)}{h}-\lim _{h \rightarrow 0} \frac{g(x+h)-g(x)}{h} \\
                & =f^{\prime}(x)-g^{\prime}(x)
            \end{align}
        \end{proof}
	
	\begin{refproof}[Teorem \ref{teorem1}]
		Referanslı ispat örneği
	\end{refproof}
	
	\subsection{Birinci bölüm üçüncü derece başlık}
	\begin{lemma}
		Örnek yardımcı teorem ifadesi
	\end{lemma}
		\subsubsection{Birinci bölüm dördüncü derece başlık}
	\begin{proposition}
		Örnek önerme ifadesi
	\end{proposition}
	
	\begin{definition}
		Örnek tanım ifadesi
	\end{definition}
	
	\begin{corollary}
		Örnek sonuc ifadesi
	\end{corollary}
	
	\begin{example}
		Örnek ifadesi
	\end{example}
	
	\begin{remark}
		Örnek not ifadesi
	\end{remark}
	
	
	Örnek tablo ve figür kodu aşağıda bulunabilir. 
	
	\begin{example}[Örnek Tablo]
		Burada dikkat edilmesi gereken önemli nokta
		\begin{verbatim}
			\caption{ ... }
		\end{verbatim}
		içerisine yazılacak olan açıklamanın aynı zamanda 
		\begin{verbatim}
			\textmd{ ... }
		\end{verbatim}
		içerisine yazılması gerekmektedir. Aksi halde açıklama yazısı kalın olacaktır. 
		\begin{table}[h]
			\caption{\textmd{$\vert (G_{{w}}^{B_{3}}g)(x)-g(x) \vert$ at some
					random values}}\centering
			\begin{tabular}{c c c}
				\hline
				\textbf{$w$} & \textbf{$\vert (G_{{w}}^{B_{3}}g)(2, 1.55)-g(2, 1.55) \vert$} & 
				\textbf{$\vert (G_{{w}}^{B_{3}}g)(0.2, 4.15)-g(0.2, 4.15) \vert$} \\ \hline
				$ 5 $ & $ 0,0959681 $ & $ 0,00403352 $  \\ 
				$ 100 $ & $ 0,000243682 $ & $ 0,0000112008 $  \\ 
				$ 500 $ & $9.74806\times 10^{-6}$ & $ 4.48069\times 10^{-7}  $ \\ 
				$ 1000 $ & $ 2.43702\times 10^{-6} $ & $ 1.12018\times 10^{-7} $ \\ \hline
			\end{tabular}
			\label{table-g}
		\end{table}
	\end{example}
	
	\begin{example}
		Burada dikkat edilmesi gereken önemli nokta \begin{verbatim} \shorthandoff{=} \end{verbatim} komutunun kullanılmasıdır. \begin{verbatim} \includegraphics \end{verbatim} komutunda özelleştirme yapılıyorsa (köşeli parantez içinde özel istekler belirtildiyse) ve komut silinirse figür hata verecektir. "=" işaretini kullanabilmemiz için yukarıdaki komutu ilgili komutun hemen önüne eklemeliyiz. Kullanım bittikten sonra (yada figürün sonunda) = işareti \begin{verbatim} \shorthandon{=} \end{verbatim} komutu ile eski haline getirilebilir.
  
		\begin{figure}
			\shorthandoff{=}
			\includegraphics[scale=1]{img/T1}
			\renewcommand{\baselinestretch}{0.5}
			\caption{\textmd{Selçuk Üniversitesi Fen Bilimleri Enstitüsü Logo}}
			\label{logo}
                \shorthandon{=}
		\end{figure}
	\end{example}
	
	
	
	\chapter{KAYNAK ARAŞTIRMASI}
	
	Kaynak araştırması bölümünü yazmaya buradan başlayınız.
	
	\section{İkinci Bölüm İkinci Derece Başlık}
	
	Paragraf içinde kullanıma göre farklı stilde atıf yapılabilir. Doğru atıf stilini kullanmak için tez yazım kuralları dosyasını inceleyiniz. Kullanılabilecek atıf komutları:
	
	\begin{enumerate}
		\item \cite{Acar2016}: Yazar (Yıl)
		\item \citealp{Acar2016}: Yazar, Yıl
		\item \citep{Acar2016}: (Yazar, Yıl)
		\item \citet*{Acar2016}: Yazar (Yıl)
		\item \citeauthor{Acar2016}: Yazar
		\item \citeyear{Acar2016}: Yıl
	\end{enumerate}
	Lütfen güncel tez şablonundan metin içinde kullanılmasına izin verilen stili kullanınız.
	
	Makale örneği: \cite{Acar2016}
	
	Kitap Örneği: \cite{Aral2013}
	
	Basılmış Tez Örneği (Yüksek Lisans): \cite{Turgay2021}
	
	Basılmış Tez Örneği (Doktora): \cite{Acar2015}
	
	
	\subsection{İkinci bölüm üçüncü derece başlık}

	\subsubsection{İkinci bölüm dördüncü derece başlık}
	
	\chapter{MATERYAL VE YÖNTEM}
	
	\section{Üçüncü Bölüm İkinci Derece Başlık}
	
	\subsection{Üçüncü bölüm üçüncü derece başlık}
	
	\subsubsection{Üçüncü bölüm dördüncü derece başlık}

	
	\chapter{ARAŞTIRMA SONUÇLARI VE TARTIŞMA}
	
	Araştırma sonuçları ve tartışma bölümünü yazmaya buradan başlayınız.
	
	\section{Dördüncü Bölüm İkinci Derece Başlık}
	
	\subsection{Dördüncü bölüm üçüncü derece başlık}
	
	\subsubsection{Dördüncü bölüm dördüncü derece başlık}

	\chapter{SONUÇLAR VE ÖNERİLER}
	
	\section{Sonuçlar}
	Sonuçlar bölümünü yazmaya buradan başlayınız.
	
	\section{Öneriler}
	Öneriler bölümünü yazmaya buradan başlayınız.\footnote{Bu bir dipnot kullanım örneğidir!}
	
	
	
	
	
	\clearpage

	\bibliographystyle{apaselcuk}
	\bibliography{tez_library}

	%%%
	%%% Özgeçmiş sayfasını oluştur.
	%%%
	\begin{ozgecmis}
		\begin{spacing}{1}
			\textbf{K\.I\c{S}\.ISEL B\.ILG\.ILER} 
			\bigskip \par
			\textbf{Ad{\i} Soyad\i :}  \par  
			\textbf{Uyru\u{g}u :}   \par
			\textbf{Do\u{g}um Yeri ve Tarihi :}  \par
			\textbf{e-mail :} 
			\bigskip \par
			\textbf{EĞ\.IT\.IM}
			\medskip \par
			\begin{tabular}{ p{5cm} p{5cm} p{5cm} }
				\textbf{Derece} & \textbf{Kurum} & \textbf{Bitirme Y\i l\i } \\
				Lise: & ... & ...\\
				Üniversite: & ... & ...\\
				Yüksek Lisans: & ... & ...
			\end{tabular}
			
			\bigskip \par
			\textbf{\.I\c{S} DENEY\.IMLER\.I}
			\medskip \par
			\begin{tabular}{ p{5cm} p{5cm} }
				\textbf{Kurum} & \textbf{Yıl}  \\
				Kurum 1: & ...\\
				Kurum 2: & ...
			\end{tabular}
		
			
			\bigskip \par
			\textbf{UZMANLIK ALANI} 
			\medskip \par
                \begin{tabular}{ p{5cm} }
				\textbf{Alan} \\
				Alan 1 \\
				Alan 2 
			\end{tabular}
			
			\bigskip \par
			\textbf{YABANCI D\.ILLER}
			\medskip \par
			\begin{tabular}{ p{5cm} p{5cm} }
				\textbf{Yabancı Dil} & \textbf{Seviye}  \\
				Yabancı Dil 1: & Seviye 1\\
				Yabancı Dil 2: & Seviye 2
			\end{tabular}
			
			\bigskip \par
			\textbf{SERTİFİKALAR}
			\medskip \par
                \begin{tabular}{ p{5cm} p{5cm} p{5cm} }
				\textbf{Sertifika Alanı} & \textbf{Verildiği Kurum} & \textbf{Yıl } \\
				Sertifika 1 & ... & ...\\
				Sertifika 2 & ... & ...
			\end{tabular}
			
			\bigskip \par
			\textbf{BEL\.IRTMEK \.ISTED\.IĞ\.IN\.IZ D\.IĞER ÖZELL\.IKLER}
			\medskip \par

							
			\bigskip \par
			\textbf{YAYINLAR}  
			\medskip \par
                Yayın Bilgisi 1 \par
                Yayın Bilgisi 2 \par
			
			\bigskip \par
			\textbf{BİLDİRİLER}  
			\medskip \par
			Bildiri Bilgisi 1 \par
                Bildiri Bilgisi 2 \par

		\end{spacing}                                  
	\end{ozgecmis} 
\end{document}                       %% Done.

